\section{Der invariante Ellipsenlift}

Es seien $(c_x,c_y)\in\R^2$ der Mittelpunkt und $a,b\in\R$ die Halbachsen.
Für einen Kontrollwert $g\in[-1,1]$ und einen Zweigwert
$s\in\{-1,1\}$ definieren wir
\begin{equation}\label{eq:lift}
  L(g,s)=
  \left(
    c_x+ag,
    c_y+bs\sqrt{1-g^2}
  \right).
\end{equation}

Statt der divisionsbehafteten Standardform verwenden wir die polynomiale
Ellipseninvariante
\begin{equation}\label{eq:polynomial-invariant}
  b^2(x-c_x)^2+a^2(y-c_y)^2=a^2b^2.
\end{equation}
Diese Form bleibt auch bei entarteten Halbachsen sinnvoll.

\subsection*{Invarianzsatz}

Für jedes $g$ mit $g^2\leq1$ und jedes $s$ mit $s^2=1$ erfüllt
$L(g,s)$ die Gleichung~\eqref{eq:polynomial-invariant}.

\subsection*{Beweis}

Einsetzen von~\eqref{eq:lift} ergibt
\begin{align*}
 &b^2(ag)^2+a^2\left(bs\sqrt{1-g^2}\right)^2 \\
 &\quad=a^2b^2\left(g^2+s^2(1-g^2)\right)\\
 &\quad=a^2b^2.
\end{align*}
Der Beweis benutzt ausschließlich $g^2\leq1$ und $s^2=1$. Insbesondere ist
unerheblich, ob $g$ durch eine lineare Funktion, eine Fourierreihe oder den
Kosinus erzeugt wurde. Genau darin liegt die Invariante des Modells.
